% Guía de consulta del endpoint genero-cliente (Bases de Datos, UP)
% Compilar con: pdflatex guia-consulta-endpoint.tex   (dos veces, o latexmk -pdf)
\documentclass[11pt,a4paper]{article}
\usepackage[utf8]{inputenc}
\usepackage[T1]{fontenc}
\usepackage[spanish,es-noshorthands]{babel}
\usepackage[margin=2.2cm]{geometry}
\usepackage{lmodern}
\usepackage{booktabs}
\usepackage{array}
\usepackage{xcolor}
\usepackage{listings}
\usepackage{hyperref}
\hypersetup{colorlinks=true, urlcolor=blue!60!black, linkcolor=black}

\newcolumntype{L}[1]{>{\raggedright\arraybackslash}p{#1}}

\definecolor{codebg}{RGB}{16,24,40}
\definecolor{codefg}{RGB}{226,232,240}
\lstset{
  basicstyle=\ttfamily\small\color{codefg},
  backgroundcolor=\color{codebg},
  frame=single, framerule=0pt, framesep=6pt, xleftmargin=6pt, xrightmargin=6pt,
  breaklines=true, breakatwhitespace=false, columns=fullflexible, keepspaces=true,
  showstringspaces=false,
  literate={ã}{{\~a}}1 {í}{{\'i}}1 {é}{{\'e}}1 {á}{{\'a}}1 {ó}{{\'o}}1 {ú}{{\'u}}1 {ñ}{{\~n}}1
}

\title{Cómo consultar el endpoint \texttt{genero-cliente}}
\author{Bases de Datos, UP}
\date{Octubre de 2026}

\begin{document}
\maketitle

Modelo que predice el género musical probable de un cliente a partir de su proveedor de correo, país y
ciudad. Entrenado cada día por un job de Databricks y publicado en un endpoint de Databricks Model Serving.

\section*{Datos de acceso}
\begin{tabular}{@{}L{3.2cm}L{12.3cm}@{}}
\toprule
\textbf{Dato} & \textbf{Valor} \tabularnewline
\midrule
URL del endpoint & \url{https://dbc-c19856c3-faa2.cloud.databricks.com/serving-endpoints/genero-cliente/invocations} \tabularnewline
Token de acceso & se entrega por separado; va en la cabecera \texttt{Authorization: Bearer <token>} \tabularnewline
Método & \texttt{POST}, cuerpo JSON, cabecera \texttt{Content-Type: application/json} \tabularnewline
\bottomrule
\end{tabular}

\section*{Entrada}
Una lista de registros con tres campos de texto:

\begin{tabular}{@{}L{2.6cm}L{8.4cm}L{4.5cm}@{}}
\toprule
\textbf{Campo} & \textbf{Qué es} & \textbf{Ejemplo} \tabularnewline
\midrule
\texttt{tipo\_correo} & proveedor del correo (\texttt{gmail}, \texttt{yahoo}, \texttt{hotmail}\ldots); también se acepta el correo completo & \texttt{yahoo} o \texttt{ana@yahoo.com.br} \tabularnewline
\texttt{pais} & país del cliente & \texttt{Germany} \tabularnewline
\texttt{ciudad} & ciudad del cliente & \texttt{Berlin} \tabularnewline
\bottomrule
\end{tabular}

\section*{Ejemplo con curl}
Sustituya \texttt{PEGUE\_AQUI\_EL\_TOKEN} por el token, sin comillas ni signos \texttt{<} \texttt{>}:

\begin{lstlisting}
curl -X POST https://dbc-c19856c3-faa2.cloud.databricks.com/serving-endpoints/genero-cliente/invocations \
  -H "Authorization: Bearer PEGUE_AQUI_EL_TOKEN" \
  -H "Content-Type: application/json" \
  -d '{"dataframe_records": [{"tipo_correo": "yahoo", "pais": "Germany", "ciudad": "Berlin"}]}'
\end{lstlisting}

En Git Bash o PowerShell de Windows, el cuerpo va con comillas dobles escapadas:

\begin{lstlisting}
curl -X POST https://dbc-c19856c3-faa2.cloud.databricks.com/serving-endpoints/genero-cliente/invocations -H "Authorization: Bearer PEGUE_AQUI_EL_TOKEN" -H "Content-Type: application/json" -d "{\"dataframe_records\": [{\"tipo_correo\": \"yahoo\", \"pais\": \"Germany\", \"ciudad\": \"Berlin\"}]}"
\end{lstlisting}

\section*{Ejemplo con Python}
\begin{lstlisting}[language=Python]
import os, requests

URL = "https://dbc-c19856c3-faa2.cloud.databricks.com/serving-endpoints/genero-cliente/invocations"
TOKEN = os.environ["DATABRICKS_TOKEN"]   # el token, en una variable de entorno

respuesta = requests.post(
    URL,
    headers={"Authorization": f"Bearer {TOKEN}", "Content-Type": "application/json"},
    json={"dataframe_records": [
        {"tipo_correo": "yahoo", "pais": "Germany", "ciudad": "Berlin"},
        {"tipo_correo": "gmail", "pais": "Brazil", "ciudad": "São Paulo"},
    ]},
    timeout=60,
)
print(respuesta.status_code)
print(respuesta.json())
\end{lstlisting}

\section*{Respuesta}
Un objeto \texttt{predictions} con un resultado por registro enviado, en el mismo orden:

\begin{lstlisting}
{
  "predictions": [
    {
      "genero": "Metal",
      "top_1": "Metal", "p_1": 0.4076,
      "top_2": "Rock", "p_2": 0.3364,
      "top_3": "Alternative & Punk", "p_3": 0.0752
    }
  ]
}
\end{lstlisting}

\texttt{genero} es el género predicho; \texttt{top\_1} a \texttt{top\_3} son los tres géneros más probables con su probabilidad.

\section*{Si algo falla}
\begin{tabular}{@{}L{4.6cm}L{5.4cm}L{5.5cm}@{}}
\toprule
\textbf{Respuesta} & \textbf{Causa} & \textbf{Solución} \tabularnewline
\midrule
\texttt{401 Credential was not sent or was of an unsupported type} & el token no llegó bien (falta, con \texttt{<} \texttt{>}, o con espacios) & revise la cabecera \texttt{Authorization: Bearer <token>} \tabularnewline
\texttt{403} & el token no tiene permiso sobre este endpoint o ya venció & pida un token nuevo \tabularnewline
la primera llamada tarda más de lo normal & el endpoint se apaga cuando no se usa y arranca con la primera petición & espere y reintente; las siguientes responden en segundos \tabularnewline
\texttt{400} con mención a columnas & falta alguno de los tres campos o el nombre está mal escrito & use exactamente \texttt{tipo\_correo}, \texttt{pais}, \texttt{ciudad} \tabularnewline
\bottomrule
\end{tabular}

\end{document}
